%!TEX root = ../template.tex
%%%%%%%%%%%%%%%%%%%%%%%%%%%%%%%%%%%%%%%%%%%%%%%%%%%%%%%%%%%%%%%%%%%%
%% chapter5.tex
%% NOVA thesis document file
%%
%% Chapter with Validation and Testing
%%%%%%%%%%%%%%%%%%%%%%%%%%%%%%%%%%%%%%%%%%%%%%%%%%%%%%%%%%%%%%%%%%%%


\typeout{NT FILE chapter5.tex}%

\chapter{Project Development}
\label{cha:proj_development}

\section{Arduino and EMG Placement and Configuration}
\label{sec:arduino_emg_placement}

\subsubsection{EMG Positioning}
\label{sssec:emg_positioning}

The muscle chosen to be measured was the Gastrocnemius Lateralis as it is one of the most important muscles for a jumping motion.
The position for the \gls{EMG} sensors was chosen as recommended by the \gls{EU}'s SENIAM project (Surface ElectroMyoGraphy for the Non-Invasive Assessment of Muscles)~\cite{emgPosition}.

%precisa de foto do posicionamento


\section{Movement Recognition - Data Processing ML Model}
\label{sec:data_processing_ML}

\section{Movement Recognition - Data Processing}
\label{sec:data_processing_VT}

